%%%%%%%%%%%%%%%%%
% This is an sample CV template created using altacv.cls
% (v1.1.4, 27 July 2018) written by LianTze Lim (liantze@gmail.com). Now compiles with pdfLaTeX, XeLaTeX and LuaLaTeX.
% 
%% It may be distributed and/or modified under the
%% conditions of the LaTeX Project Public License, either version 1.3
%% of this license or (at your option) any later version.
%% The latest version of this license is in
%%    http://www.latex-project.org/lppl.txt
%% and version 1.3 or later is part of all distributions of LaTeX
%% version 2003/12/01 or later.
%%%%%%%%%%%%%%%%

%% If you need to pass whatever options to xcolor
\PassOptionsToPackage{dvipsnames}{xcolor}

%% If you are using \orcid or academicons
%% icons, make sure you have the academicons 
%% option here, and compile with XeLaTeX
%% or LuaLaTeX.
% \documentclass[10pt,a4paper,academicons]{altacv}

%% Use the "normalphoto" option if you want a normal photo instead of cropped to a circle
% \documentclass[10pt,a4paper,normalphoto]{altacv}

\documentclass[10pt,a4paper,UTF8]{altacv}
%% AltaCV uses the fontawesome and academicon fonts
%% and packages. 
%% See texdoc.net/pkg/fontawecome and http://texdoc.net/pkg/academicons for full list of symbols.
%% 
%% Compile with LuaLaTeX for best results. If you
%% want to use XeLaTeX, you may need to install
%% Academicons.ttf in your operating system's font 
%% folder.


% Change the page layout if you need to
\geometry{left=1cm,right=9cm,marginparwidth=6.8cm,marginparsep=1.2cm,top=1.25cm,bottom=1.25cm,footskip=2\baselineskip}

% Change the font if you want to.

% If using pdflatex:
\usepackage[T1]{fontenc}
\usepackage[utf8]{inputenc}
\usepackage[default]{lato}
\usepackage[UTF8]{ctex}
% If using xelatex or lualatex:
% \setmainfont{Lato}

% Change the colours if you want to
\definecolor{Mulberry}{HTML}{72243D}
\definecolor{SlateGrey}{HTML}{2E2E2E}
\definecolor{LightGrey}{HTML}{666666}
\colorlet{heading}{Sepia}
\colorlet{accent}{Mulberry}
\colorlet{emphasis}{SlateGrey}
\colorlet{body}{LightGrey}

% Change the bullets for itemize and rating marker
% for \cvskill if you want to
\renewcommand{\itemmarker}{{\small\textbullet}}
\renewcommand{\ratingmarker}{\faCircle}
%% sample.bib contains your publications
\addbibresource{sample.bib}

\usepackage[colorlinks]{hyperref}

\begin{document}

\name{苏向宇}
% \tagline{求职意向：算法实习生}
\photo{3.5cm}{1.png}
\personalinfo{%
  % Not all of these are required!
  % You can add your own with \printinfo{symbol}{detail}
  % \email{xiangyv.su@gmail.com }
  \phone{+86-15112314526}
  \email{xiangyv.su@gmail.com }\\
  
%   \github{github.com/XiangyuSu611}
  %\linkedin{linkedin.com/in/GargiBhave}
  \mailaddress{深圳大学可视计算研究中心}
  \homepage{\href{https://xiangyusu611.github.io/}{\textcolor{black}{xiangyusu611.github.io}}}
  %\twitter{@twitterhandle}
  %\linkedin{linkedin.com/in/yourid}
  
  %% You MUST add the academicons option to \documentclass, then compile with LuaLaTeX or XeLaTeX, if you want to use \orcid or other academicons commands.
%   \orcid{orcid.org/0000-0000-0000-0000}
}

%% Make the header extend all the way to the right, if you want. 
\begin{fullwidth}
\makecvheader
\end{fullwidth}

%% Depending on your tastes, you may want to make fonts of itemize environments slightly smaller
% \AtBeginEnvironment{itemize}{\small}


%% Provide the file name containing the sidebar contents as an optional parameter to \cvsection.
%% You can always just use \marginpar{...} if you do
%% not need to align the top of the contents to any
%% \cvsection title in the "main" bar.
\cvsection[page1sidebar]{教育经历}

\cvevent{\textbf{计算机科学与技术（博士即将入学）- 深圳大学}}{}{2023年9月始}{广东深圳}


\divider

\cvevent{\textbf{计算机技术（硕士学位）- 深圳大学}}{}{2020年9月 - 2023年7月}{广东深圳}
\begin{itemize}
% \item GPA: 3.43/4（专业排名第九名）
\item 2020-2022年均获深圳大学研究生学业奖学金一等奖
\end{itemize}

\divider

\cvevent{\textbf{计算机科学与技术（学士学位）- 沈阳建筑大学}}{}{2015年9月 - 2019年7月}{辽宁沈阳}
\begin{itemize}
\item 2015-2019年均获沈阳建筑大学优秀学生奖学金
\end{itemize}

\cvsection{项目经历}

\cvevent{\textbf{基于材质重建的三维模型材质迁移\\Reconstruction-based Material Transfer}}{}{2022年6月 - 至今}{}
\begin{itemize}
\item 该项目旨在解决现有的材质迁移方法对输入数据的限制，在二维图像和三维模型之间建立更加灵活的对应关系，并直接从图像中重建真实感材质，实现高质量材质迁移。
\end{itemize}

\divider

\cvevent{\textbf{基于图像的三维模型跨结构材质迁移\\Photo-to-Shape Material Transfer for Diverse Structures}}{ \footnotesize{\textcolor[rgb]{0.6,0.6,0.6}{Ruizhen Hu, Xiangyu Su, Xiangkai Chen, Oliver van Kaick, Hui Huang}}   \\ ACM Transactions on Graphics (SIGGRAPH 2022) }{2021年3月 - 2022年3月}{}
\begin{itemize}
\item https://vcc.tech/research/2022/TMT
\item 该项目以真实照片作为指导，自动为三维模型添加高质量真实感材质，可自动生成大量带有高质量材质的三维模型。
\end{itemize}

\cvsection{实习经历}
\cvevent{\textbf{腾讯AI Lab 视觉计算组研究实习生}}{}{2022年8月 - 2023年2月}{}


\medskip

\clearpage




%% If the NEXT page doesn't start with a \cvsection but you'd
%% still like to add a sidebar, then use this command on THIS
%% page to add it. The optional argument lets you pull up the 
%% sidebar a bit so that it looks aligned with the top of the
%% main column.
% \addnextpagesidebar[-1ex]{page3sidebar}

\end{document}
